\documentclass[platex,dvipdfmx,12pt]{jsarticle}
%使ってるのがplatexだからクラス指定はplatex.間違えてもuplatexにしないように

%==============================================================================
% 出欠管理システム 要件定義書（版 0.8 ドラフト）
%   原稿 : 要件定義書_レビュー用.html
%   書式 : requirements-template.tex（kosen-report-writer のプリアンブル + 要件定義書サンプル.pdf の体裁）
%
% ■ビルド
%   latexmk -outdir=latex-output requirements-attendance.tex
%   （.latexmkrc により platex → dvipdfmx で処理される）
% ■図は assets/ に置いた PNG を \reqimage で読み込む（無ければ枠が出る）
%
% ■独自マクロ一覧（プリアンブルで定義）
%   \reqlabel{名前}       ■見出し。直後に \begin{...}（表）が続くときは表の先頭行に組み込まれ，
%                         見出しだけがページ末に取り残されない（段落が続くときは単独の見出し）
%   \reqth{文字}          表ヘッダー用（白・太字）
%   \reqimage[幅]{パス}    画像．無ければ枠表示
%   reqbullet 環境        本文の「・」箇条書き（\item で項目）
%   reqtable 環境         ヘッダー付き表   \begin{reqtable}{列指定}{ヘッダー行}
%   reqlist 環境          ヘッダー無しの2列表（項目名＋値）\begin{reqlist}{列指定}
%   reqcategory 環境      カテゴリー | 詳細 の表（1章）
%   reqitems 環境         項目 | 説明 | 具体的な要件 の表（4章）
%   列指定 : K{比率}=太字ゴシックの項目列  P{比率}=通常列  M{比率}=中央揃え列
%            比率は表幅に対する割合で，合計が 1 になるようにする（例 K{.3}P{.7}）
%==============================================================================

\usepackage{graphicx}
\usepackage[hyphens]{url}
\usepackage{otf}
\usepackage{multirow}
\usepackage{amsmath}
\usepackage{here}
\usepackage{geometry}
\usepackage{ascmac}

\usepackage[hang,small,bf]{caption}
\usepackage[subrefformat=parens]{subcaption}
\captionsetup{compatibility=false}
\usepackage{textcomp}

\usepackage{listings} %日本語のコメントアウトをする場合jvlisting（もしくはjlisting）が必要
%ここからソースコードの表示に関する設定
\lstset{
  basicstyle={\ttfamily},              % 基本:等幅フォント
  identifierstyle={\small},            % 識別子
  commentstyle={\smallitshape},        % コメント:斜体
  keywordstyle={\small\bfseries},     % キーワード:太字
  stringstyle={\small\ttfamily},      % 文字列
  frame={tb},                           % 上部と下部に線を表示
  breaklines=true,                      % 行が長い場合に折り返す
  columns=[l]{fullflexible},            % 列幅を自動調整する（見た目が良くなる）
  numbers=left,                         % 行番号を左側に表示
  xrightmargin=0zw,                     % 右マージンのサイズ．
  xleftmargin=1.6zw,                    % 左マージンのサイズ．行番号が2桁でも行左端からはみ出ない値．
  numberstyle={\scriptsize},           % 行番号のスタイル．スクリプトサイズのフォントを使用
  stepnumber=1,                         % 何行ごとに行番号を表示するか
  numbersep=1zw,                        % 行番号とソースコードの間の距離
  lineskip=-1.4ex,                      % ソースコードの行間（結構詰め気味）
  morestring=[b]",                     %日本語等のStringを正常にダブルクオーテーションで囲む
  showstringspaces=false,               %Stringに囲まれた部分ではスペースがSymbol付きで表示されるのでそれを消す
}
\renewcommand{\lstlistingname}{ソースコード}

\geometry{left=20mm,right=20mm,top=20mm,bottom=30mm}
\renewcommand{\baselinestretch}{1.2}

%------------------------------------------------------------------------------
% ここから要件定義書用の追加（上までがスキル標準のプリアンブル）
%   longtable : ページをまたぐ長い表（ヘッダー行を各ページに再掲）
%   needspace : 見出し・■見出しだけがページ末尾に取り残されるのを防ぐ
%   xcolor[table] : サンプルの水色ヘッダー・縞模様・青枠のための色指定
%------------------------------------------------------------------------------
\usepackage{longtable}
\usepackage{needspace}
\usepackage[table]{xcolor}

% サンプルPDFの配色
\definecolor{reqblue}{RGB}{68,114,196}   % 表紙の枠・見出しの青
\definecolor{reqbox}{RGB}{217,226,243}   % 章見出しの塗り
\definecolor{reqhead}{RGB}{91,155,213}   % 表ヘッダー
\definecolor{reqlight}{RGB}{222,234,246} % 表の縞（薄い水色）
\definecolor{reqrule}{RGB}{157,195,230}  % 表の罫線

\makeatletter %sectionをchapterっぽくする.
% 見出し番号を「1.」「(1)」形式にする（サンプルPDFに合わせた設定）．
% スキル標準の「第1章」形式にする場合は pre=第 / post=章 に戻し，\thesubsection の行を削除する
\renewcommand{\presectionname}{}
\renewcommand{\postsectionname}{.}%はコメントですがつけ忘れると\section*でバグります
\renewcommand{\thesubsection}{(\arabic{subsection})}

% 図番号を1.5のようにする
\renewcommand{\theequation}{\arabic{section}.\arabic{equation}}
\renewcommand{\thefigure}{\arabic{section}.\arabic{figure}}
\renewcommand{\thetable}{\arabic{section}.\arabic{table}}
% jsarticle は章ごとにカウンタがリセットされないため追加（無いと2章の最初の図が「図2.2」になる）
\@addtoreset{equation}{section}
\@addtoreset{figure}{section}
\@addtoreset{table}{section}

% 章見出し : 青枠＋薄青塗りのボックス（スタイル末尾のマクロが見出し全体を受け取る）
\newcommand{\reqsectionbox}[1]{%
  \begingroup\setlength{\fboxrule}{1.6pt}\setlength{\fboxsep}{5pt}%
  \noindent\fcolorbox{reqblue}{reqbox}{\parbox{\dimexpr\linewidth-2\fboxsep-2\fboxrule\relax}{\raggedright #1}}%
  \endgroup}
% 節見出し : 左に青いバー，下に青い罫線
\newcommand{\reqsubsectionbox}[1]{%
  \noindent\textcolor{reqblue}{\rule[-.8ex]{.5zw}{3ex}}\hspace{.5zw}%
  \begin{minipage}[t]{\dimexpr\linewidth-1zw\relax}\raggedright #1\end{minipage}\par\nobreak
  \vspace{-.7\baselineskip}\noindent\textcolor{reqblue}{\rule{\linewidth}{1.2pt}}\par}

\renewcommand{\section}{%
\clearpage
\@startsection{section}{1}{\z@}%
{\Cvs \@plus.5\Cdp \@minus.2\Cdp}% 見出し前に付けられる空白量
{.5\Cvs \@plus.3\Cdp}% 見出し後に付けられる垂直方向の空白量
{\normalfont\Large\headfont\reqsectionbox}}% 見出しのスタイル
\renewcommand{\subsection}{%
\Needspace{6\baselineskip}%
\@startsection{subsection}{2}{\z@}%
{\Cvs \@plus.5\Cdp \@minus.2\Cdp}%
{.5\Cvs \@plus.3\Cdp}%
{\normalfont\large\headfont\reqsubsectionbox}}
\makeatother

\renewcommand{\contentsname}{目次}
\setcounter{tocdepth}{2}

\makeatletter
% 表の列数を数える（列指定 K{..}P{..}M{..} の文字数）
\newcount\req@ncols
\newcommand{\req@countcols}[1]{\req@ncols\z@
  \@tfor\req@t:=#1\do{%
    \def\req@a{K}\ifx\req@t\req@a\advance\req@ncols\@ne\fi
    \def\req@a{P}\ifx\req@t\req@a\advance\req@ncols\@ne\fi
    \def\req@a{M}\ifx\req@t\req@a\advance\req@ncols\@ne\fi}}

% ■見出し
%   直後が \begin{...}（表）のときは見出しを表の先頭行に組み込む（見出しだけがページ末に取り残されない）。
%   表が続かないとき（段落が続くとき）は単独の見出しとして出す。
\def\req@label{}
\newcommand{\req@plainlabel}[1]{\par\Needspace{4\baselineskip}\medskip\noindent ■#1\par\nopagebreak\smallskip}
\newcommand{\reqlabel}[1]{\@ifnextchar\begin{\gdef\req@label{#1}}{\req@plainlabel{#1}}}
\newcommand{\req@labelrow}{\ifx\req@label\@empty\else
  \multicolumn{\the\req@ncols}{@{}l@{}}{■\req@label}\\[.6ex]\fi}

% 本文の「・」箇条書き
\newenvironment{reqbullet}{%
  \begin{list}{・}{\setlength{\leftmargin}{1.5zw}\setlength{\labelwidth}{1zw}\setlength{\labelsep}{.5zw}%
    \setlength{\itemsep}{0pt}\setlength{\parsep}{0pt}\setlength{\topsep}{.3\baselineskip}}%
}{\end{list}}

% 画像（ファイルが無い間はテンプレート単体でコンパイルが通るよう枠を表示する）
\newcommand{\reqimage}[2][.9\textwidth]{%
  \IfFileExists{#2}%
    {\includegraphics[width=#1]{#2}}%
    {\fbox{\parbox[c][5cm][c]{\dimexpr#1-2\fboxsep-2\fboxrule\relax}{\centering ※図を添付（\detokenize{#2}）}}}}

% 表 : 列指定。比率は表幅に対する割合（セル余白と罫線ぶんを差し引いて表幅ぴったりに収める）
%   K・P列は左揃え（\raggedright）。狭い列で両端揃えにすると字間が不自然に空くため
\newcommand{\reqth}[1]{\textcolor{white}{\bfseries #1}}
\newcolumntype{K}[1]{|>{\bfseries\gtfamily\raggedright\arraybackslash}p{\dimexpr#1\dimexpr\linewidth-\arrayrulewidth\relax-2\tabcolsep-\arrayrulewidth\relax}}
\newcolumntype{P}[1]{|>{\raggedright\arraybackslash}p{\dimexpr#1\dimexpr\linewidth-\arrayrulewidth\relax-2\tabcolsep-\arrayrulewidth\relax}}
\newcolumntype{M}[1]{|>{\centering\arraybackslash}p{\dimexpr#1\dimexpr\linewidth-\arrayrulewidth\relax-2\tabcolsep-\arrayrulewidth\relax}}

% ヘッダー付き表（最初のデータ行が薄水色の縞模様）。直前の \reqlabel は先頭行として組み込まれる
\newenvironment{reqtable}[2]{%
  \par\smallskip\small
  \renewcommand{\arraystretch}{1.25}%
  \req@countcols{#1}%
  \arrayrulecolor{reqrule}%
  \ifx\req@label\@empty\rowcolors{2}{reqlight}{white}\else\rowcolors{2}{white}{reqlight}\fi% 先頭用・継続用ヘッダーの2行ぶんを数えたうえで，最初のデータ行が水色になる向き
  \begin{longtable}{#1|}
  \req@labelrow
  \hline
  \rowcolor{reqhead}#2\\
  \hline
  \endfirsthead
  \hline
  \rowcolor{reqhead}#2\\
  \hline
  \endhead
}{%
  \end{longtable}\gdef\req@label{}\par\smallskip}
% ヘッダー無しの表（項目名＋値）
\newenvironment{reqlist}[1]{%
  \par\smallskip\small
  \renewcommand{\arraystretch}{1.25}%
  \req@countcols{#1}%
  \arrayrulecolor{reqrule}%
  \ifx\req@label\@empty\rowcolors{1}{reqlight}{white}\else\rowcolors{2}{white}{reqlight}\fi
  \begin{longtable}{#1|}
  \req@labelrow
  \hline
  \endfirsthead
  \hline
  \endhead
}{%
  \end{longtable}\gdef\req@label{}\par\smallskip}
% よく使う表のプリセット
\newenvironment{reqcategory}{%
  \begin{reqtable}{K{.3}P{.7}}{\reqth{カテゴリー} & \reqth{詳細}}}{\end{reqtable}}
\newenvironment{reqitems}{%
  \begin{reqtable}{K{.25}P{.38}P{.37}}{\reqth{項目} & \reqth{説明} & \reqth{具体的な要件}}}{\end{reqtable}}
\makeatother


\begin{document}

\begin{titlepage}
\begin{center}
\vspace*{6em}
{\LARGE －版 0.8（ドラフト）－}

\vspace{6em}
\begingroup\setlength{\fboxrule}{3pt}\setlength{\fboxsep}{0pt}%
\fcolorbox{reqblue}{white}{\parbox[c][.26\textheight][c]{.72\textwidth}{\centering
{\huge\headfont 出欠管理システム}\\[.6em]
{\huge\headfont 要件定義書}}}
\endgroup

\vspace{9em}
{\Large 2026年10月19日}

\vspace{1.5em}
{\Large ShuMASUI}
\end{center}
\end{titlepage}
\setcounter{page}{2}

\section*{改訂履歴}
\begin{reqtable}{P{.14}P{.18}P{.16}P{.52}}{\reqth{バージョン} & \reqth{日付} & \reqth{改訂者} & \reqth{改訂内容}}
0.1 & 2026-10-16 & ShuMASUI & 議事録をもとに新規作成 \\ \hline
0.2 & 2026-10-16 & ShuMASUI & 要件定義書サンプルの構成に合わせて再構成（本タブ） \\ \hline
0.3 & 2026-10-16 & ShuMASUI & ヒアリングの回答を反映（授業の構造、規模、性能目標、QRコード方式、テナント構成など） \\ \hline
0.4 & 2026-10-16 & ShuMASUI & ドメインモデル図と回答を反映（公欠の登録、任意の修正、アカウントの付与、AWS、バックアップ、冗長化など） \\ \hline
0.5 & 2026-10-16 & ShuMASUI & 回答を反映（出席状態の区分、欠課時数、開始・終了時限、公欠の扱い、任意の申請への統合、スケジュールなど） \\ \hline
0.6 & 2026-10-19 & ShuMASUI & QRコード方式の出席確認フローを反映（管理者によるQRコードの発行、所属の検査、認証の検証、位置情報の照合、代返への対策） \\ \hline
0.7 & 2026-10-19 & ShuMASUI & 回答を反映（授業の回数、先生が持つ授業の数） \\ \hline
0.8 & 2026-10-19 & ShuMASUI & 回答を反映（設計時に判断する事項、認証の移行方式、DBのリカバリ、位置情報、QRコードの再発行、学外からの利用、テスト、移行、運用体制など）。引継ぎと教育の項を削除 \\ \hline
\end{reqtable}

\tableofcontents

\section{概要}
本書は、奈良高専の先生へのヒアリング議事録と、これまでの要件定義方針をもとに、出欠管理システムの要件をまとめたものである。

\subsection{システム化の目的}
\begin{reqbullet}
  \item 授業ごと・先生ごとにバラバラな出席管理を、統一システムに集約する。
  \item 先生と学生が同じ出席情報を見て、記録に誤りがないかをチェックできるようにする。
  \item 学校へ提出する最終確定の出席情報を、CSVで出力できるようにする。
  \item 現状の運用を壊さずに導入する。
\end{reqbullet}

\subsection{現状の課題}
\begin{reqtable}{K{.3}P{.7}}{\reqth{カテゴリー} & \reqth{課題}}
出席管理 & ・出席の管理方法が、授業ごと・先生ごとにバラバラである。\par・授業ごとの出席状況は、先生が個人で管理している。\par・先生と学生が、記録に誤りがないかを統一システムで確認できない。\par・遅刻よりも、欠課時数の処理が問題になる。 \\ \hline
学校への提出 & ・最終確定の出席情報を、学校側のExcelに入力している。 \\ \hline
登録・運用 & ・学生を一人ずつ登録する作業が面倒で、運用の負担になっている。 \\ \hline
出席確認の操作 & ・一般教科では、スマホを出してボタンを押す操作を面倒と感じる先生が多い。\par・WebClassを使っていない先生もいる。 \\ \hline
\end{reqtable}

\subsection{プロジェクトの背景}
\begin{reqcategory}
現場の利用状況 & ・専門科目・演習科目では、WebClassを使う前提の授業がある。一般教科では使っていない先生もいる。\par・Teamsは、コロナ禍に使い始めた先生が多い。資料の配布や学生への連絡に使われている。\par・授業の開始時に学生にスマホを使わせることは、先生が認めれば規則上は問題ない。 \\ \hline
導入上の制約 & ・現状の運用を壊せないことが、ヒアリングで最初に挙がった課題感である。\par・WebClassやTeamsからは独立したアプリケーションとし、内部リンクは一切貼らない。\par・Microsoft SSOは申請に時間がかかるため、MVPでは使わない。認証はCognitoに委任し、学校が付与したMicrosoftアカウントをメインキーとして保存して、サービスの利用を始める。 \\ \hline
開発方針 & ・出席確認型のマルチテナントSaaSとし、AWS上の完全サーバーレス環境で構築する。学校テナントの下に、先生単位の小テナントを置く。要件はドメインモデルv1.1をベースに定義する。\par・将来は、Microsoft SSOに対応し、本人確認を兼ねて、より堅牢にする。対応後も、再登録なしでユーザーデータを使い続けられるようにする。 \\ \hline
運用の前提 & ・学生は、公欠をシステムに登録する。公欠は、システムでは判別できない。奈良高専の教務課が管理し、突合して対応する。従来のフローと変わらない。\par・先生は、何も登録しない。割り当てられた授業の、割り当てられた内容に対してのみ操作する。\par・オペレータと運用・保守担当は、ベンダー側の担当である。アカウントは、ベンダーが全ユーザーへ付与する。シードアカウントは用意しない。 \\ \hline
\end{reqcategory}

\subsection{システム構成図}
\begin{figure}[H]
  \begin{center}
    \reqimage[\textwidth]{assets/system-architecture.png}
    \caption{システム構成図：利用者・端末・システム・外部の関係}
    \label{fig:system-architecture}
  \end{center}
\end{figure}

先生・学生・オペレータが出欠管理システムを使い、学生の出席確認では、管理者が発行したQRコードを学生が読み取り、位置情報とあわせて送る。認証はCognitoに委任し、学校へはCSVで出席情報を渡す。

\subsection{用語定義}
\begin{reqtable}{K{.3}P{.7}}{\reqth{用語} & \reqth{意味}}
出欠管理システム & 本システム。授業での出席確認、出席状況の確認、学校提出用のCSV出力を行うSaaS。 \\ \hline
出席確認 & 授業の開始時に、学生が出席したことをシステムに記録すること。QRコードの読み取りと、位置情報の照合で確認する。 \\ \hline
出席状態 & 出席、欠席、早退、遅刻、遅刻早退、公欠の6つ。公欠は、システムでは判別できない。 \\ \hline
欠課時数 & 授業を欠課した時間の数。欠席、遅刻、早退などから処理する。 \\ \hline
先生 & 授業を担当し、出席状況を確認する利用者。何も登録せず、割り当てられた授業の、割り当てられた内容に対してのみ操作する。 \\ \hline
学生 & 授業に登録され、出席確認を行う利用者。 \\ \hline
オペレータ & ベンダー側の担当者。利用者（学生・先生・オペレータ）と授業を、システムに登録する。 \\ \hline
運用・保守担当 & ベンダー側で、システムの運用と保守を行い、アカウントを付与する担当者。 \\ \hline
マルチテナント & 一つのシステムを複数の組織で使い、データを組織ごとに分離する方式。 \\ \hline
テナント & マルチテナントでデータを分離する単位。学校テナントと、先生単位の小テナントがある。学校テナントは、ベンダーが追加する。 \\ \hline
小テナント & 学校テナントの下に置く、先生単位のテナント。 \\ \hline
メンバーシップ & 学生を、先生単位のテナントに結びつける関係。 \\ \hline
Bastion & 外部から内部の環境へ安全に接続するための、中継用のサーバー。ベンダーが、学校テナントの追加に使う。 \\ \hline
Amazon Cognito & 本システムが認証を委任するAWSのサービス。 \\ \hline
SSO & 一度のログインで複数のサービスを使える仕組み。将来、Microsoft SSOに対応する。 \\ \hline
多要素認証 & パスワードに加えて、別の確認方法を使う認証。本システムでは用意せず、SSOの統合後に高専のポリシーで自動的に適用される。 \\ \hline
メインキー & 同一人物かどうかを判定するための情報。学校が付与したMicrosoftアカウントを使う。 \\ \hline
前期・後期 & 年度を前半と後半に分けた期間。 \\ \hline
半期授業 & 前期のみ、または後期のみで行う授業。 \\ \hline
通年授業 & 前期と後期を通して行う授業。 \\ \hline
時限 & 授業の時間の区切り。1限から8限があり、時間の概念として持つ。授業は、開始時限と終了時限を持つ。 \\ \hline
公欠 & 学校が公に認める欠席。状態としては持つが、システムでは判別できず、奈良高専の教務課が処理する。 \\ \hline
教務課 & 奈良高専の部署。公欠を判別し、管理し、突合する。 \\ \hline
任意の申請 & 出席結果の誤りや、端末を忘れた場合などに、学生が出す申請。先生が承認する。 \\ \hline
WebClass & 学校の授業で使われている学習管理システム。 \\ \hline
Teams & Microsoft Teams。授業資料の配布や学生への連絡に使われている。 \\ \hline
MVP & 最初にリリースする、必要最小限の機能の範囲。 \\ \hline
CSV & カンマ区切りのデータ形式。学校へ提出する出席情報の出力形式として使う。 \\ \hline
管理者 & 授業の出席確認用QRコードを発行する利用者。先生を想定しているが、確認が必要。 \\ \hline
QRコード & 出席確認ページのURLを含むコード。URLのクエリが授業を示す。管理者が発行し、学生が読み取る。 \\ \hline
位置情報 & 学生がQRコードを読み取ったときに、端末から一緒に送られる場所の情報。学校の中で読み取ったかの照合のみに使い、保存しない。 \\ \hline
代返 & 他人による代理の出席確認。所属の検査と、位置情報の照合で防ぐ。 \\ \hline
\end{reqtable}

\section{業務要件}

\subsection{業務要件一覧}
本システムは、以下の業務を実現するものとする。出席確認の方式は、QRコードの読み取りと、位置情報の照合とする。流れは、(2) 業務フロー図に示す。

\reqlabel{出席確認}
\begin{reqtable}{K{.2}P{.42}P{.22}M{.16}}{\reqth{業務} & \reqth{詳細} & \reqth{実施者} & \reqth{システム化対象}}
出席確認 & 授業の開始時に、管理者が発行した授業のQRコードを学生が読み取り、出席確認を行う。ログインした学生がその授業に所属しているかと、読み取ったときの位置情報が学校の中かを、システムが検証する。 & 学生・管理者（先生を想定） & ○ \\ \hline
QRコードの発行 & 管理者が、授業の出席確認用QRコードを発行する。 & 管理者（先生を想定） & ○ \\ \hline
出席確認の検証 & 所属の検査、認証の検証、位置情報の照合を行い、学校の中で読み取ったことを確認する。代返への対策にもなる。 & システム & ○ \\ \hline
出席状況の確認 & 先生は割り当てられた授業の出席状況を、学生は自分の出席状況を確認する。 & 先生・学生 & ○ \\ \hline
欠課時数の処理 & 欠席・早退・遅刻などの出席状態から、欠課時数を処理できるようにする。算出の規則は【未定】。 & システム & ○ \\ \hline
公欠の登録 & 学生が、公欠をシステムに登録する。公欠かどうかは、システムでは判別できない。 & 学生 & ○ \\ \hline
任意の申請と承認 & 学生が、出席結果の誤りや端末の忘れなどを、任意の申請として出す。先生が承認すると、申請された状態で、そのユーザーのデータを更新する。 & 学生・先生 & ○ \\ \hline
出席状態の任意の修正 & 自らが登録されている授業の情報に限り、出席状態を任意で変更する。 & ユーザー（役割は要確認） & ○ \\ \hline
\end{reqtable}

\reqlabel{登録}
\begin{reqtable}{K{.2}P{.42}P{.22}M{.16}}{\reqth{業務} & \reqth{詳細} & \reqth{実施者} & \reqth{システム化対象}}
利用者・授業の登録 & 名簿をまとめてシステムに登録する。先生は何も登録しない。ベンダー側の処理も入るため、資料を受け取ってから7営業日以内に行う。 & オペレータ & ○ \\ \hline
ログイン情報の案内 & 登録した利用者へ、ログインに必要な情報をメールで送る。 & システム & ○ \\ \hline
\end{reqtable}

\reqlabel{学校への提出}
\begin{reqtable}{K{.2}P{.42}P{.22}M{.16}}{\reqth{業務} & \reqth{詳細} & \reqth{実施者} & \reqth{システム化対象}}
出席情報のCSV出力 & 最終確定した出席情報を、CSVで出力する。形式は、最後まで未定の見込み。 & 先生 & ○ \\ \hline
学校側Excelへの入力 & CSVの情報を、学校側のExcelに入力する。実施者は要確認。 & 先生 & × \\ \hline
公欠の判別・管理と突合 & 公欠の判別と管理は、奈良高専の教務課が行い、突合して対応する。従来のフローと変わらない。 & 教務課 & × \\ \hline
\end{reqtable}

\reqlabel{運用・管理}
\begin{reqtable}{K{.2}P{.42}P{.22}M{.16}}{\reqth{業務} & \reqth{詳細} & \reqth{実施者} & \reqth{システム化対象}}
利用者の認証 & ログインの認証を、Cognitoに委任する。 & システム & ○ \\ \hline
テナント・アカウントの発行 & ベンダーが、学校テナントを追加し、全ユーザーへアカウントを付与する。シードアカウントは用意しない。 & ベンダー & × \\ \hline
\end{reqtable}

\subsection{業務フロー図}
\begin{figure}[H]
  \begin{center}
    \reqimage[.85\textwidth]{assets/business-flow.png}
    \caption{業務フロー図：登録から学校への提出まで、5つの担当で示す}
    \label{fig:business-flow}
  \end{center}
\end{figure}

オペレータの一括登録から始まり、学生が授業の開始時に、QRコードを読み取って出席確認を行う（流れは下表）。誤りがなければ、先生が確認して最終確定する。誤りがあれば、学生が申請し、先生が承認する。最終確定した情報はCSVで出力し、学校側のExcelへの入力は現状どおり行う。公欠は、学生がシステムに登録する。教務課は、従来どおり公欠を管理し、突合して対応する。

\reqlabel{出席確認（QRコード）のフロー}
\begin{reqtable}{M{.08}P{.24}P{.68}}{\reqth{手順} & \reqth{実施者} & \reqth{内容}}
1 & 管理者（先生を想定） & 授業の出席確認用QRコードを、授業ごとに発行しなおす。QRコードにURLを含め、URLのクエリが授業を示す。 \\ \hline
2 & 学生 & 端末のカメラアプリなどでQRコードを読み取り、URLのページへ移動する。 \\ \hline
3 & システム & クエリを照合して、どの授業かを特定する。 \\ \hline
4 & システム & 学生のログイン情報で、その授業に本当に所属しているかを検査する。 \\ \hline
5 & システム & アカウントの認証ができているかを、セッションなどで検証する。検証の方式は【未定】。 \\ \hline
6 & 学生の端末 & 3〜5の検証を通った状態で読み取っていれば、読み取りの情報を、位置情報とあわせて送る。 \\ \hline
7 & システム & 「QRコードを読み取れた事実」と、「同時に送られた位置情報」を照合し、本当に学校の中で読み取ったかを確認する。 \\ \hline
8 & システム & 照合に通れば、出席確認として記録する。 \\ \hline
\end{reqtable}

この流れは、他人による代理の出席確認（代返）への対策も兼ねる。位置情報の照合の細部と、QRコードの使い回しの扱いは、末尾の確認事項に挙げた。

\subsection{規模}
\begin{reqlist}{K{.3}P{.7}}
先生数 & 約100名 \\ \hline
学生数 & 1000名を超える \\ \hline
1授業の学生数 & 約40名 \\ \hline
先生と授業の関係 & 先生1人が複数の授業を担当する。大半の先生は、1〜5程度の授業を持つ。先生は何も登録せず、割り当てられた授業の、割り当てられた内容に対してのみ操作する。先生単位の小テナントの下に授業があり、学生はメンバーシップで結びつく。 \\ \hline
授業の期間 & 年度は、前期と後期に分かれる。前期のみ、後期のみの半期授業と、前期後期にわたる通年授業がある。 \\ \hline
授業の時限 & 1限から8限まである。時限は、時間の概念として持つ。授業は、開始時限と終了時限を持つ。多くの科目は、隣り合う2限分で行う。 \\ \hline
授業の回数 & 1つの開講された講座について、前期16回、後期16回。通年授業は32回、半期授業は16回で、授業の開講期間に応じて変動する。 \\ \hline
テナント数 & 学校テナントの数は【未定】。小テナントは、先生の数（約100）に応じる。 \\ \hline
出席確認の同時利用 & 授業の開始時に、全学年・全学科が同時にアクセスする。40名×5クラス×7学年=1400リクエストを想定する。 \\ \hline
データ件数 & 最大のオーダーは、授業の数×学生の数。 \\ \hline
\end{reqlist}

\subsection{場所}
\begin{reqbullet}
  \item 出席確認: 教室で行う。学生が、端末のカメラアプリなどでQRコードを読み取る。位置情報で、学校の中にいることを確認する。
  \item 出席状況の確認・申請・承認・登録: 学校のネットワークの外からも使える。
\end{reqbullet}

\subsection{時期・時間}
\begin{reqlist}{K{.3}P{.7}}
出席確認の時間 & 授業が実施されている時間のみ（1限から8限） \\ \hline
アクセスのピーク & 授業の開始時 \\ \hline
システム稼働時間 & 24時間365日。出席確認は授業の時間のみだが、確認などの作業を含め、すべての操作が使える。 \\ \hline
メンテナンス & 半年に1回のメジャーアップデートで、1時間ほど停止する恐れがある。緊急の障害対応で、停止する場合もある。 \\ \hline
学校への提出の時期 & 概ね7月から8月と、2月から3月 \\ \hline
\end{reqlist}

\subsection{指標}
数値目標は、ヒアリングで示された値を記載する。決まっていない項目は【未定】とした。

\reqlabel{性能}
\begin{reqtable}{P{.5}P{.5}}{\reqth{指標} & \reqth{目標}}
授業の開始時に、全学年・全学科が同時にアクセスしても処理できる。 & 40名×5クラス×7学年=1400リクエストを、1秒以内に処理する \\ \hline
学校へ提出する出席情報を、CSVで出力する。 & 10秒以内 \\ \hline
\end{reqtable}

\reqlabel{操作性}
\begin{reqtable}{P{.5}P{.5}}{\reqth{指標} & \reqth{目標}}
先生の操作負担を減らす。現状は2回以上の操作が必要である。 & 先生は何も登録しない。1度の登録のみで、以降は全自動で管理される \\ \hline
\end{reqtable}

\reqlabel{品質}
\begin{reqtable}{P{.5}P{.5}}{\reqth{指標} & \reqth{目標}}
システムが24時間365日使える。 & 稼働率はイレブンナイン（99.999999999\%）を目指す【要確認】 \\ \hline
障害から復旧するまでの時間。 & MTTR（平均復旧時間）60分を目指す \\ \hline
欠課時数を処理できる。 & 算出の規則は【未定】 \\ \hline
画面の品質。 & Lighthouseの、パフォーマンス、アクセシビリティ、ベストプラクティスの3項目すべてで、90点を超える \\ \hline
\end{reqtable}

\reqlabel{セキュリティ}
\begin{reqtable}{P{.5}P{.5}}{\reqth{指標} & \reqth{目標}}
システムによる情報漏洩を防ぐ。 & 漏洩0件。マルチテナントを分離する \\ \hline
\end{reqtable}

\subsection{システム化範囲}
\reqlabel{システム化の対象範囲}
\begin{reqtable}{K{.3}P{.7}}{\reqth{業務} & \reqth{内容}}
出席確認 & 出席確認（QRコードの発行と読み取り、所属の検査、認証の検証、位置情報の照合）、出席状況の確認、欠課時数の処理、公欠の登録、任意の申請と承認、出席状態の任意の修正 \\ \hline
登録 & 利用者・授業の一括登録、ログイン情報のメール案内 \\ \hline
学校への提出 & 出席情報のCSV出力 \\ \hline
運用・管理 & 利用者の認証（Cognito）、マルチテナント対応 \\ \hline
\end{reqtable}

\reqlabel{システム化の対象外範囲}
\begin{reqtable}{K{.3}P{.7}}{\reqth{業務} & \reqth{内容}}
学校への提出 & 学校側Excelへの入力（学校側の作業）、学校側システムとの直接連携（CSV出力で代替する）、公欠の判別・管理と突合（教務課が従来どおり行う。システムでは判別できない） \\ \hline
運用・管理 & 学校テナントの追加と、全ユーザーへのアカウントの付与（ベンダーの作業。シードアカウントは用意しない） \\ \hline
連携 & WebClassやTeamsとの連携（独立したアプリケーションとし、内部リンクも貼らない） \\ \hline
認証 & 多要素認証（本システムでは用意しない。SSOの統合後は、高専のポリシーで自動的に適用される） \\ \hline
将来対応 & Microsoft SSOへの対応（再登録なしで、ユーザーデータを使い続けられるようにする） \\ \hline
\end{reqtable}

\section{機能要件}

\subsection{機能一覧}
\reqlabel{基本機能}
\begin{reqtable}{M{.10}K{.16}P{.36}P{.16}P{.22}}{\reqth{ID} & \reqth{機能} & \reqth{説明} & \reqth{操作主体} & \reqth{対応業務要件}}
F-01 & 出席確認 & 学生が授業の開始時に、管理者が発行したQRコードを読み取って、出席確認を行える。読み取りでは、授業の特定、所属の検査、認証の検証、位置情報の照合を行う（F-13、F-14）。 & 学生・管理者（先生を想定） & 出席確認 \\ \hline
F-02 & 出席状況の確認 & 先生は割り当てられた授業の出席状況を、学生は自分の出席状況を確認できる。 & 先生・学生 & 出席状況の確認 \\ \hline
F-03 & 利用者の一括登録 & オペレータが、利用者（学生・先生・オペレータ）の情報をまとめて登録できる。氏名や、学校が付与したMicrosoftアカウントなどを登録する。1件ずつ紐付ける手作業を前提にしない。先生は何も登録しない。 & オペレータ & 利用者・授業の登録 \\ \hline
F-04 & 授業の登録 & 授業、担当の先生、授業に登録される学生の関係を登録できる。授業には、年度、期（前期・後期）、半期か通年か、開始時限と終了時限を持たせる。 & オペレータ & 利用者・授業の登録 \\ \hline
F-05 & ログイン情報の案内 & 登録した利用者へ、ログインに必要な情報をメールで自動的に送る。 & システム & ログイン情報の案内 \\ \hline
F-06 & CSV出力 & 学校へ提出する最終確定の出席情報を、CSVで出力できる。形式は、最後まで未定の見込み。 & 先生 & 出席情報のCSV出力 \\ \hline
F-07 & 任意の申請と承認（出席結果の誤りなど） & 学生が、出席結果の誤りや端末の忘れなどを、任意の申請として統一的に出す。先生は、その申請を承認するのみとする。承認した場合は、申請された状態で、そのユーザーのデータを更新する。 & 学生・先生 & 任意の申請と承認 \\ \hline
F-08 & ログイン・認証 & 利用者がログインできる。認証はCognitoが行う。学校が付与したMicrosoftアカウントをメインキーとして保存し、それに紐づけてサービスの利用を始める。 & 先生・学生・オペレータ & 利用者の認証 \\ \hline
F-09 & メンバーシップの管理 & 学生を、先生単位のテナントにメンバーシップで結びつけられる。 & オペレータ & 利用者・授業の登録 \\ \hline
F-10 & 公欠の登録 & 学生が、公欠をシステムに登録できる。公欠かどうかは、システムでは判別できない。判別と処理は、奈良高専の教務課が行う。 & 学生 & 公欠の登録 \\ \hline
F-11 & 出席状態の任意の修正 & 自らが登録されている授業の情報に限り、出席状態を任意で変更できる。障害などで出席確認ができない場合も、この機能で各ユーザーが対応する。 & ユーザー（役割は要確認） & 出席状態の任意の修正 \\ \hline
F-12 & 欠課時数の処理 & 欠席・早退・遅刻・遅刻早退などの出席状態から、欠課時数を処理できる。算出の規則は【未定】。 & システム & 欠課時数の処理 \\ \hline
F-13 & QRコードの発行 & 管理者が、授業の出席確認用QRコードを発行できる。QRコードには出席確認ページのURLを含め、URLのクエリで授業を示す。授業ごとに発行しなおす。有効期間は【未定】。 & 管理者（先生を想定） & 出席確認 \\ \hline
F-14 & 出席確認の検証 & QRコードを読み取った学生について、次の順で検証する。(1) クエリから、どの授業かを特定する。(2) ログイン情報で、その授業に所属しているかを検査する。(3) 認証ができているかを、セッションなどで検証する。(4) 読み取りの事実と、同時に送られた位置情報を照合し、学校の中で読み取ったかを確認する。すべて通れば、出席確認として記録する。 & システム & 出席確認 \\ \hline
\end{reqtable}

\reqlabel{拡張機能（将来対応）}
\begin{reqtable}{M{.10}K{.2}P{.45}P{.25}}{\reqth{ID} & \reqth{機能} & \reqth{説明} & \reqth{対応業務要件}}
E-01 & Microsoft SSOへの対応 & 将来、Microsoft SSOに対応する。対応した時点で、SSOの情報を突合し、登録済みのユーザーデータを使い続けられるようにする（再登録の手間をなくす）。方式は、CognitoのIdentity Provider連携と、IDでの突合とする。 & 利用者の認証 \\ \hline
\end{reqtable}

\reqlabel{出席確認方式}
出席確認は、QRコードを読み取る方式とする。ICカード方式は採用しない。流れは、業務要件の(2) 業務フロー図の表のとおり。管理者がQRコードを発行し、学生が読み取ったときに、ログインの情報と位置情報で検証する。

\subsection{画面}
画面は、機能一覧から導いた案である。PC版とスマートフォン版の2種類を用意する。ブラウザは、Chrome、Safari、Firefoxの最新版を対象にする。

\reqlabel{共通}
\begin{reqtable}{K{.3}P{.7}}{\reqth{画面名} & \reqth{説明}}
ログイン画面 & Cognitoの認証でログインする。 \\ \hline
\end{reqtable}

\reqlabel{学生向け}
\begin{reqtable}{K{.3}P{.7}}{\reqth{画面名} & \reqth{説明}}
出席確認画面 & カメラアプリなどでQRコードを読み取ると、URLの出席確認ページが開く。ログイン状態の確認と、位置情報の送信の許可を経て、出席確認を行う。 \\ \hline
出席状況確認画面 & 自分の出席状況を、授業ごとに確認する。 \\ \hline
公欠登録画面 & 公欠を登録する。 \\ \hline
申請画面 & 出席結果の誤りや端末の忘れなどを、任意の申請として出す。 \\ \hline
\end{reqtable}

\reqlabel{先生向け}
\begin{reqtable}{K{.3}P{.7}}{\reqth{画面名} & \reqth{説明}}
QRコード発行画面 & 管理者が、授業の出席確認用QRコードを発行し、表示する。 \\ \hline
出席状況確認画面 & 割り当てられた授業の出席状況を、学生ごとに確認する。 \\ \hline
申請承認画面 & 学生からの申請を確認し、承認する。承認すると、申請された状態で、そのユーザーのデータを更新する。 \\ \hline
CSV出力画面 & 学校へ提出する出席情報を、CSVで出力する。 \\ \hline
\end{reqtable}

\reqlabel{要確認}
\begin{reqtable}{K{.3}P{.7}}{\reqth{画面名} & \reqth{説明}}
出席状態の修正画面 & 自らが登録されている授業の情報に限り、出席状態を任意で変更する。使える利用者の役割は要確認。 \\ \hline
\end{reqtable}

\reqlabel{オペレータ向け}
\begin{reqtable}{K{.3}P{.7}}{\reqth{画面名} & \reqth{説明}}
利用者一括登録画面 & 学生・先生・オペレータの名簿をまとめて登録する。 \\ \hline
授業登録画面 & 授業、担当の先生、授業に登録される学生の関係を登録する。 \\ \hline
\end{reqtable}

\subsection{帳票}
\begin{reqtable}{K{.28}P{.54}P{.18}}{\reqth{帳票名} & \reqth{内容} & \reqth{形式}}
学校提出用の出席情報 & 最終確定した出席情報。項目と列構成は、最後まで未定のままになる見込み。 & CSV \\ \hline
\end{reqtable}

その他の帳票は【未定】である。

\subsection{情報・データ}
利用者とシステムのやり取りは、下図のドメインモデルを想定する。ER図は、別途添付する【要追記】。

\begin{figure}[H]
  \begin{center}
    \reqimage[\textwidth]{assets/domain-model.png}
    \caption{ドメインモデル：学校側とベンダー側の、本システムとのやり取り}
    \label{fig:domain-model}
  \end{center}
\end{figure}

先生と学生は、本人の証明（ログイン）と出欠確認を行い、出欠状況を確認する。学生は、公欠も登録する。ベンダー側では、運用・保守担当がアカウントを付与し、オペレータが人と授業を登録する。

\begin{reqtable}{K{.26}P{.74}}{\reqth{情報・データ名} & \reqth{内容}}
テナント情報 & 学校テナントと、その下の先生単位の小テナント。学校テナントは、ベンダーがBastion経由で専用ツールを用いて追加する。 \\ \hline
利用者情報 & 先生・学生・オペレータの氏名、学校が付与したMicrosoftアカウント、役割。Microsoftアカウントは、メインキーとして保存し、将来のSSOで同一人物を突合する。 \\ \hline
メンバーシップ & 学生と、先生単位のテナントを結びつける情報。 \\ \hline
授業情報 & 授業、担当の先生、授業に登録された学生の関係。年度、期（前期・後期）、半期か通年か、開始時限と終了時限を持つ。時限は、1限から8限の時間の概念として持つ。 \\ \hline
出席記録 & 授業の回ごとの、学生の出席状態。状態は、出席、欠席、早退、遅刻、遅刻早退、公欠の6つ。公欠は、システムでは判別できない。件数の最大のオーダーは、授業の数×学生の数になる。 \\ \hline
欠課時数 & 出席状態から処理する、欠課の時間の数。算出の規則は【未定】。 \\ \hline
公欠登録 & 学生が登録した公欠の内容。 \\ \hline
任意の申請 & 学生が出した申請の内容（出席結果の誤り、端末の忘れなど）と、先生の承認結果。承認すると、申請された状態で、そのユーザーのデータを更新する。 \\ \hline
出席状態の修正記録 & 任意の申請を承認する経路と、任意で修正する経路の、両方の修正の記録（案）。 \\ \hline
学校提出用データ & 最終確定した出席情報。CSVで出力する。形式は、最後まで未定の見込み。 \\ \hline
出席確認用QRコード & 管理者が発行する、授業のQRコード。出席確認ページのURLを含み、URLのクエリが授業を示す。授業ごとに発行しなおす。有効期間は【未定】。 \\ \hline
位置情報 & 出席確認のときに、学生の端末から送られる。学校の中で読み取ったかの照合のみに使い、保存しない。IPアドレスのログのみ取る。 \\ \hline
\end{reqtable}

\subsection{外部インターフェース}
\begin{reqtable}{K{.2}P{.3}P{.2}P{.3}}{\reqth{項目} & \reqth{説明} & \reqth{必要性} & \reqth{データ交換の仕様}}
認証・認可（Amazon Cognito） & 利用者の認証をCognitoに委任する。学校が付与したMicrosoftアカウントを、メインキーとして保存する。 & Microsoft SSOの申請に時間がかかるMVPで、認証をシンプルに実現するため。 & Cognitoの標準の方式。詳細は設計時に判断する。 \\ \hline
メール送信 & ログイン情報を、登録した利用者へメールで送る。 & 運用側の手作業を減らすため。 & 送信方式（サービス）は、設計時に判断する。 \\ \hline
学校（CSV） & 最終確定の出席情報を、CSVで学校へ渡す。 & 学校への提出が引き続き必要なため。 & CSVファイル。列構成と文字コードは、最後まで未定のままになる見込み。 \\ \hline
Microsoft SSO（将来） & 将来、SSOでログインできるようにする。 & より堅牢なシステムにするため。SSOの情報を突合し、再登録なしでユーザーデータを使い続けるため。 & CognitoのIdentity Provider連携を使い、IDで突合する。 \\ \hline
出席確認（QRコード） & 管理者が発行したQRコードを、学生が読み取って、出席確認を行う。QRコードにはURLを含め、URLのクエリが授業を示す。読み取りは、端末のカメラアプリなどで行う。 & 学生の出席確認を受け付けるため。 & QRコードの内容（URLとクエリの形式）は、設計時に決める【未定】。 \\ \hline
WebClass・Teams & 連携しない。独立したアプリケーションとし、内部リンクも貼らない。 & WebClassやTeamsに依存しないため。 & なし \\ \hline
教務課（公欠） & システムとは連携しない。学生が公欠をシステムに登録する。公欠は、システムでは判別できない。教務課が従来どおり、公欠を判別・管理し、突合する。 & 従来のフローを変えないため。 & なし \\ \hline
暗号化通信 & システムと利用者の端末の間の通信を暗号化する。 & 個人情報を守るため。 & TLS 1.3を使う。 \\ \hline
\end{reqtable}

\section{非機能要件}

ヒアリングで示された目標を記載し、決まっていない項目は【未定】とした。（案）は、チームの案で未合意の項目である。

\subsection{ユーザビリティおよびアクセシビリティ}
\reqlabel{ユーザビリティ要件}
\begin{reqitems}
操作の簡単さ & 先生の操作負担を減らす。 & 現状は2回以上の操作が必要である。先生は何も登録せず、1度の登録のみで、以降は全自動で管理される状態を目指す。 \\ \hline
登録の手間 & 利用者の登録の手間を減らす。 & 利用者を1件ずつ登録せず、名簿でまとめて登録できる。ベンダー側の処理も入るため、資料を受け取ってから7営業日以内に登録する。 \\ \hline
エラーメッセージ & 誤りが起きたとき、次の操作が分かるようにする。 & 原因と次の操作を示すメッセージを表示する（案）。 \\ \hline
利用する端末 & PC版とスマートフォン版の2種類を用意する。 & 対象のブラウザは、Chrome、Safari、Firefoxの最新版。 \\ \hline
\end{reqitems}

\reqlabel{アクセシビリティ要件}
\begin{reqitems}
準拠する基準 & WCAG 2.2に準拠する。 & Lighthouseのスコアは、パフォーマンス、アクセシビリティ、ベストプラクティスの3項目すべてで、90点を超える。SEOは、Webアプリケーションのため対象外とする。 \\ \hline
表示言語 & 日本語と英語のみに対応する。 & 他の言語には対応しない。 \\ \hline
\end{reqitems}

\subsection{システム方式}
\begin{reqitems}
実行基盤 & AWS上の、完全クラウドのサーバーレス環境で構築する。 & 認証にはAmazon Cognitoを使う。 \\ \hline
提供形態 & 出席確認型のマルチテナントSaaSとする。 & 学校テナントの下に、先生単位の小テナントを置く。学生はメンバーシップで結びつける。学校テナントは、ベンダーがBastion経由で専用ツールを用いて追加する。 \\ \hline
認証 & 認証をCognitoに全面的に委任する。 & 学校が付与したMicrosoftアカウントをメインキーとして保存し、それに紐づけてサービスの利用を始める。 \\ \hline
データベース & 利用者・授業・出席記録を保存する。 & DynamoDBまたはAurora Serverlessのいずれかになる見込み。どちらにするかは、設計時に判断する。 \\ \hline
監視・ログ & 監査ログは、CloudWatchへ出力されるログをベースに検証する。 & CloudWatchへ出力されるログを使う。 \\ \hline
\end{reqitems}

\subsection{規模}
業務要件の「規模」に従う。

\begin{reqitems}
利用者数 & 先生は約100名、学生は1000名を超える。 & 1授業あたりの学生は約40名。将来は、要相談のうえ、数万のオーダーまで対応できるようにする。 \\ \hline
最大同時接続ユーザー数 & 授業の開始時に、全学年・全学科が同時にアクセスする。 & 40名×5クラス×7学年=1400リクエストが同時に発生する。 \\ \hline
日次アクティブユーザー数 & 出席確認でアクセスする以外の用途について。 & 100名程度 \\ \hline
総データ量 & システム全体で管理するデータ量。 & 件数の最大のオーダーは、授業の数×学生の数。総量は未決のままとする。 \\ \hline
トランザクション数 & 1日・ピーク時に処理する件数。 & ピーク時は1400リクエスト。1日の件数は【未定】。 \\ \hline
\end{reqitems}

\subsection{性能}
\begin{reqitems}
出席確認の応答時間 & 授業の開始時に、全学年・全学科が同時にアクセスしても、学生を待たせずに記録する。 & 40名×5クラス×7学年=1400リクエストを、1秒以内に処理する。 \\ \hline
一括登録の所要時間 & 利用者の名簿をまとめて登録する。ベンダー側の処理も入る。 & 資料を受け取ってから、7営業日以内。 \\ \hline
CSV出力の処理時間 & 出席情報をCSVで出力する。 & 10秒以内。 \\ \hline
\end{reqitems}

\subsection{信頼性}
\begin{reqitems}
可用性（稼働率） & 確認などの作業を含め、すべての操作が24時間365日使える。 & 稼働率の目標はイレブンナイン（99.999999999\%）。意味と前提は【要確認】。 \\ \hline
障害対応 & 障害が起きたときの復旧の目標。 & MTTR（平均復旧時間）は60分を目指す。 \\ \hline
端末を忘れた場合 & 学生が端末を忘れた場合の扱い。 & 学生が、任意の申請を出す。先生が承認すると、申請された状態で、そのユーザーのデータを更新する。 \\ \hline
障害で出席確認ができない場合 & 障害などで出席確認ができない場合の代替方法。 & 自らが登録されている授業の情報に限り、出席状態を任意で変更できる。各ユーザーが自分で対応できる。システム全体が止まった場合の扱いは【未定】。 \\ \hline
データの耐久性 & 出席記録は学校への提出の根拠になるため、失われないようにする。 & 永続化の単一障害点は、データベースのみ。1か月に1回、データベース全体のバックアップを取り、12件分を残し続ける。1か月に1回で足りる。 \\ \hline
修正の記録 & 出席結果の修正は、記録が残る（案）。 & 任意の申請を承認する経路と、任意で修正できる経路の、2つを残す。保存の期間は【未定】。 \\ \hline
\end{reqitems}

\subsection{拡張性}
\begin{reqitems}
テナントの追加 & 新しい学校テナントを追加できる。 & ベンダー側が、Bastion経由で専用ツールを用いて追加する。 \\ \hline
認証方式の移行 & 将来、Microsoft SSOに対応する。 & 対応した時点で、SSOの情報を突合し、登録済みのユーザーデータを使い続けられるようにする（再登録の手間をなくす）。方式は、CognitoのIdentity Provider連携と、IDでの突合とする。 \\ \hline
利用者・授業数の増加 & 利用者と授業が増えても使える。 & 要相談のうえ、数万のオーダーまで対応できるようにする。 \\ \hline
\end{reqitems}

\subsection{互換性}
\begin{reqitems}
端末・ブラウザ & PC版とスマートフォン版の2種類を用意する。 & Chrome、Safari、Firefoxの最新版を、動作の要件とする。 \\ \hline
データ形式 & 学校へCSVで出力する。 & 列構成と文字コードは、最後まで未定のままになる見込み。出力する項目と文字コードを、後から変えられるようにする（案）。 \\ \hline
他システムとの連携 & WebClassやTeamsとは連携しない。 & 独立したアプリケーションとし、内部リンクも貼らない。 \\ \hline
認証の方式 & Cognitoで認証する。将来はMicrosoft SSOに対応する。 & CognitoのIdentity Provider連携を使い、IDで突合して移行する。 \\ \hline
\end{reqitems}

\subsection{継続性}
メンテナンス、冗長化、バックアップ、BCPの方針は、次のとおり。そのほかは決まっていない。

\begin{reqtable}{K{.3}P{.7}}{\reqth{項目} & \reqth{具体的な要件}}
計画メンテナンス & 半年に1回のメジャーアップデートで、1時間ほど停止する恐れがある。 \\ \hline
緊急メンテナンス & 緊急の障害対応で、停止する場合がある。 \\ \hline
冗長化 & AWSのマルチAZ環境を使う。 \\ \hline
バックアップとリカバリ & データベースが永続化の単一障害点になるため、1か月に1回、データベース全体のバックアップを取り、12件分を残し続ける。リカバリは、リードレプリカを置き、同一のデータを同期して、データベースが常に2つ存在する状態にする（RDSの場合）。 \\ \hline
サービスレベル（SLA） & 稼働率の目標とMTTRは、信頼性の項に記載したとおり。 \\ \hline
事業継続計画（BCP） & 内部向けに作成する。 \\ \hline
\end{reqtable}

\subsection{セキュリティ}

\reqlabel{データ保護}

\begin{reqitems}
個人情報の扱い & 学生の氏名・Microsoftアカウント・出席情報・位置情報を個人情報として扱う。 & 取り扱いの方針は、学校と確認して決める【未定】。 \\ \hline
通信の暗号化 & システムと利用者の端末の間の通信を暗号化する。 & TLS 1.3を使う。 \\ \hline
情報漏洩 & システムによる漏洩を0件にする。 & マルチテナントを分離する。 \\ \hline
テナント間の分離 & 他のテナントのデータを見られないようにする。 & 学校テナントごと、先生単位の小テナントごとに、データを分離する。 \\ \hline
位置情報の扱い & 出席確認のときに、端末から送られる位置情報の扱い。 & 学校の中で読み取ったかの照合のみに使い、保存しない。IPアドレスのログのみ取る。 \\ \hline
代返への対策 & 他人による代理の出席確認を防ぐ。 & 所属の検査と、読み取りの事実と位置情報の照合で対策する。QRコードは、授業ごとに発行しなおす。 \\ \hline
\end{reqitems}

\reqlabel{認証と認可}

\begin{reqitems}
認証 & 認証済みの利用者だけが使える。 & Cognitoに委任する。多要素認証は、本システムでは用意しない。Microsoft SSOの統合後は、高専のポリシーで自動的に多要素認証が適用される。 \\ \hline
アクセス権限 & 役割ごとに、見られる情報を分ける（案）。 & 自らが登録されている授業の情報のみを扱える。学生は自分の出席状況、先生は割り当てられた授業の出席状況を見られる。オペレータは登録の機能を使う。 \\ \hline
アカウントの付与 & シードアカウントは用意しない。 & ベンダーが、オペレータを含む全ユーザーへ、アカウントを付与する。 \\ \hline
\end{reqitems}

\reqlabel{脆弱性管理}

脆弱性への対応方針は【未定】である。

\subsection{稼働環境}

\begin{reqitems}
サーバー & AWS上の、完全クラウドのサーバーレス環境で動作する。 & 具体的なサービスは、設計時に判断する。 \\ \hline
利用者の端末 & PC版とスマートフォン版を用意する。 & Chrome、Safari、Firefoxの最新版。 \\ \hline
出席確認の機器 & 学生の端末のカメラアプリなどで、QRコードを読み取る。専用の読み取り機器は使わない。 & QRコードの表示方法（先生のPC、教室の投影など）は【未定】。 \\ \hline
教室のネットワーク & 教室から出席確認の通信ができる。 & 学校のネットワークの外からも、アクセスできる。 \\ \hline
\end{reqitems}

\subsection{テスト}

\begin{reqitems}
受け入れテスト & 先生と学校の要望を満たしているか確認する。 & 発注者のチームから代表者2名を選び、ロールプレイ形式で行う。合格の基準は、代表者から問題なしの判定と署名をもらい、リリースすること。 \\ \hline
テストデータ & テストデータは、ダミーデータを生成して使う。実在の学生の個人情報は使わない。 & 実データを加工したものではなく、まったくのダミーデータにする。 \\ \hline
その他のテスト & 単体、結合、システム、E2Eのテスト。 & 単体テストはファイルごと、結合テストはモジュールごとに行う。システムテストは、バックエンドとフロントエンドを独立して行う。E2Eテストは、ブラウザ操作を含めて全自動で実行する。性能・セキュリティのテストの内容は【未定】。 \\ \hline
\end{reqitems}

\subsection{移行}

現状、出席状況は先生が個人で管理している。現状の運用を壊さずに導入する必要がある。

\begin{reqitems}
既存データの移行 & 先生が個人で管理している出席の記録。 & 既存データは移行しない。 \\ \hline
現行の運用との並行 & 現行の管理とシステムを並行して使う期間。 & 現行の運用と並行して保存できるように、最終のCSVを出力できるようにしている。並行する期間は【未定】。 \\ \hline
\end{reqitems}

\subsection{運用}

\begin{reqitems}
運用体制 & ベンダー側の運用・保守担当とオペレータが運用する。オペレータが、利用者・授業の登録を行う。 & チーム5名で対応する。 \\ \hline
アカウントの付与 & ベンダーが、全ユーザーへアカウントを付与する。シードアカウントは用意しない。 & オペレータが、まとめて登録する。学校からの申請から7営業日で対応する。 \\ \hline
テナントの追加 & ベンダーが、Bastion経由で専用ツールを用いて、学校テナントを追加する。 & 内部向けに手順を作成する。 \\ \hline
日常運用 & 利用者・授業の登録と更新。 & 一括登録は、資料を受け取ってから7営業日以内。実施の時期と頻度は【未定】。 \\ \hline
問い合わせ対応 & 先生・学生からの問い合わせを受ける。 & 窓口は【未定】。 \\ \hline
障害・インシデント & 障害が起きたときの対応の流れ。 & MTTR60分を目指す。流れは【未定】。 \\ \hline
監視 & システムの稼働状況を見る。 & 監査ログは、CloudWatchへ出力されるログをベースに検証する。そのほかは【未定】。 \\ \hline
\end{reqitems}

\subsection{保守}

\begin{reqitems}
保守体制 & ベンダーの運用・保守担当が担当する。 & 体制は【未定】。 \\ \hline
定期保守 & 半年に1回のメジャーアップデートを行う。 & 1時間ほど停止する恐れがある。時期と内容は【未定】。 \\ \hline
障害対応 & 障害時の連絡と復旧。 & 緊急の障害対応で停止する場合がある。MTTR60分を目指す。連絡窓口は【未定】。 \\ \hline
機能改善 & 先生・学校からの要望の受付。 & 優先順位の付け方は【未定】。 \\ \hline
\end{reqitems}

\section{補足}

\subsection{プロジェクト スケジュール}
スケジュールは、次のとおり。リリースは2027年2月で、開発・テストは2026年12月前半、受け入れテストは12月後半に行う。

\begin{figure}[H]
  \begin{center}
    \reqimage[\textwidth]{assets/project-schedule.png}
    \caption{スケジュール：要件定義・設計からリリースまで}
    \label{fig:project-schedule}
  \end{center}
\end{figure}

\begin{reqtable}{K{.2}P{.37}P{.19}P{.24}}{\reqth{フェーズ} & \reqth{内容} & \reqth{担当} & \reqth{期間}}
要件定義・設計 & 本書の作成と、先生・学校への確認。設計。 & 開発側・顧客側 & 2026年10月〜11月 \\ \hline
開発・テスト & 開発と、開発側のテスト & 開発側 & 2026年12月前半 \\ \hline
受け入れテスト & 先生と学校の要望を満たしているかの確認 & 顧客側 & 2026年12月後半 \\ \hline
リリース & システムの公開 & 開発側 & 2027年2月 \\ \hline
\end{reqtable}

\subsection{プロジェクト 体制}
\begin{reqtable}{K{.18}P{.32}P{.5}}{\reqth{区分} & \reqth{役割} & \reqth{担当}}
顧客側 & 要件の確認、受け入れテスト & 【未定】。ヒアリングを行った奈良高専の先生を想定する。 \\ \hline
開発側 & 責任者 & 【未定】 \\ \hline
開発側 & PM & 【未定】 \\ \hline
開発側 & メンバー & 【未定】 \\ \hline
\end{reqtable}

\section*{未定事項・確認事項}
\addcontentsline{toc}{section}{未定事項・確認事項}

以下の13点が未確定である。回答済みの項目は本文に反映した。議事録は自動文字起こしで聞き取れない箇所が多いため、本書の一部は解釈を含む。回答をもらった項目から、本文の【未定】を更新する。

\begin{reqtable}{M{.08}P{.31}P{.17}P{.44}}{\reqth{No.} & \reqth{確認事項} & \reqth{確認先} & \reqth{理由}}
1 & QRコード方式の細部。管理者とは誰か（先生か、オペレータか）、QRコードの有効期間、認証の検証の方式、読み取ったときに未ログインの場合の流れ & 先生 & 管理者が授業のQRコードを発行し、学生が読み取る流れは分かった。QRコードは授業ごとに発行しなおすと分かったが、「授業ごと」は授業を行うたびの意味と解釈した。管理者の役割、有効期間、ログイン済みかの検証の方式（議事録の「譲渡」は、トークンのことと解釈した）が読み取れない \\ \hline
2 & 授業が複数の時限にまたがる場合の出席確認の単位（授業ごとに1回か、時限ごとか） & 先生 & 欠課時数を時限単位で処理するかに影響する。出席記録の件数にも影響する（授業の回ごとに記録すると、授業の数×学生の数×回数（16〜32回）になる） \\ \hline
3 & 欠課時数の算出の規則（遅刻・早退・遅刻早退を、何時間分の欠課とみなすか） & 先生・学校 & 欠課時数を処理できるようにする要望があるが、規則が示されていない。CSVの形式は未定のままの見込みだが、欠課時数を出力に含めるかも確認したい \\ \hline
4 & 任意の申請（出席結果の誤り、端末の忘れなど）の期限と、却下の有無。申請の種類を分けるか & 先生・チーム & 「承認するのみ」とあるが、期限や却下の扱いが書かれていない \\ \hline
5 & 出席状態の任意の修正（F-11）を行える役割（先生のみか、学生も含むか） & チーム & 範囲は「自らが登録されている授業の情報のみ」と分かったが、役割が読み取れない \\ \hline
6 & 公欠の状態を、誰がどのようにつけるか & 先生・教務課 & 図では学生が公欠を登録するが、公欠は「システムでは判別できない」とある。教務課が、システムの情報をどう突合するかも確認したい \\ \hline
7 & 議事録の最後に出た依頼や、次回の打ち合わせの時期 & 先生 & 文字起こしが不明瞭で、次の行動が読み取れない \\ \hline
8 & 学校側Excelへの入力を行うのは、先生か事務か & 先生 & 議事録では実施者が読み取れない \\ \hline
9 & 稼働率の「イレブンナイン」（99.999999999\%）の意味 & チーム & 稼働率なら、年間の停止の許容は約0.3ミリ秒になり、半年に1回約1時間のメンテナンスと両立しない。データの耐久性を指すか、計画メンテナンスを除く稼働率かを確認したい \\ \hline
10 & 個人情報の取り扱い方針（学生の氏名・Microsoftアカウント・出席情報・位置情報） & 学校 & 学校の規則に従う必要がある \\ \hline
11 & プロジェクトの体制 & チーム & 情報がなく、【未定】のまま \\ \hline
12 & 位置情報の照合の細部。「学校の中」の範囲の決め方（敷地か、教室か）、屋内での精度、位置情報を許可しない・取得できない場合の扱い & 先生・学校 & 位置情報は個人情報にあたる。許可されない端末で出席確認ができない場合の扱い（任意の申請を使うか）も決める必要がある。議事録の「大変な対策」は、代返の対策と解釈した \\ \hline
13 & QRコードを読み取った時刻による、出席と遅刻の判定の基準。早退と遅刻早退の記録の方法 & 先生 & 出席状態は6つあるが、QRコードを読み取るのは授業の開始時とのことで、早退などをどう記録するかが読み取れない \\ \hline
\end{reqtable}

\begin{flushright}
以上
\end{flushright}

\end{document}
